%% kapa: exam=1 round=37 year=2026 period=2 form=A subject=appraisal_law subject_ko=감정평가관계법규 seq=8 number=8
%% source: Q-Net 감정평가사 기출문제 (cst00309) artlSeq=5258924 file=raw/1008/5258924/2026년 제37회 감정평가사 1차 2교시.pdf
%% conversion: kordoc 4.24.51 -> tools/gen1.py (auto); review=proofread-sonnet; flags=none
\begin{question}{8}
국토의 계획 및 이용에 관한 법령상 지구단위계획구역의 지정에 관한 설명으로 옳은 것은?
\begin{choices}
  \item ｢주택법｣에 따른 대지조성사업지구에 대하여는 지구단위계획구역을 지정할 수 없다.
  \item ｢택지개발촉진법｣에 따라 지정된 택지개발지구에서 시행되는 사업이 끝난 후 5년이 지난 지역은 지구단위계획구역으로 지정하여야 한다.
  \item 계획관리지역을 지구단위계획구역으로 지정하려면 지정하려는 구역 면적의 100분의 30 이상이 계획관리지역이어야 한다.
  \item 계획관리지역을 지구단위계획구역으로 지정하려면 지구단위계획구역에 계획관리지역 외에 보전녹지지역 또는 보전관리지역이 포함되어야 한다.
  \item 개발진흥지구로서 도시지역 외에 있는 지역을 지구단위계획구역으로 지정하려면 자연환경ㆍ경관ㆍ미관 등을 해치지 아니하고 국가유산의 훼손우려가 없는 지역이어야 한다.
\end{choices}
\end{question}
\answer{5}
